\section{Cold jets and uniform color estimates}
\label{jet:section}

We construct the common finite-entropy datum and bound each cold color
by its free transport plus a small Maxwellian remainder. Narrow velocity
profiles control self-collisions, while spatial gaps make collisions
between distinct colors small. These estimates supply the cold density
bound and the mixed-cold source bound used to select the hot caps in
Section~\ref{sec:cap}.

The estimates in this section hold for every fixed choice
\[
 K\ge1,\qquad 0<\eta\le\log2,\qquad
 0<\alpha\le\tfrac14,\qquad 0<\gamma\le1.
\]
These four parameters will subsequently be fixed by the growth argument.
Here we choose the geometric and thermal parameters after them.
To keep the later shortening of the time interval separate from the
estimates, fix the preliminary bounds
\begin{equation}
 \overline s=\tfrac18,\qquad \overline T=\tfrac1{32},
 \qquad 0<s_0\le\overline s,\qquad 0<T_0\le\overline T .
 \label{jet:preliminary}
\end{equation}
All estimates proved below hold on the preliminary interval
\([0,\overline T]\).  Their constants do not change when \(T_0\) is
subsequently decreased.  They may depend on the fixed parameters and on
the fixed function \(J\), but never on the number of retained annuli,
the collision cutoff, or the size of the hot component.

\subsection{The geometry and the datum}

Put
\[
 s_j=s_0e^{-j\eta},\qquad p_j=s_j^2,\qquad
 c_\eta=\frac{e^\eta-1}{16},\qquad
 \delta_j=\frac{c_\eta}{4}s_j^2,\qquad j\ge1.
\]
For \(s=s_j\), define the annular core and the two stripe patterns by
\begin{align*}
 A_j&=\{y\in\R^2:
       s_j+c_\eta p_j<|y|<e^\eta s_j-c_\eta p_j\},\\
 P_j^+&=\bigcup_{k\in\mathbb Z}
       \{y:kp_j<y_1<(k+\tfrac14)p_j\},\\
 P_j^-&=\bigcup_{k\in\mathbb Z}
       \{y:(k+\tfrac12)p_j<y_1<(k+\tfrac34)p_j\},\\
 E_j^\sigma&=A_j\cap P_j^\sigma,\qquad \sigma\in\{-1,+1\}.
\end{align*}
Each stripe pattern has period \(p_j\) and fraction
\(\theta=1/4\).  This is a statement about the periodic pattern:
intersecting it with the annulus need not preserve the fraction
exactly.  We retain the original sets \(E_j^\sigma\) in the datum.

Let
\[
 \widehat E_j^\sigma
   =\{y:\dist(y,E_j^\sigma)\le\delta_j\},
 \qquad a=\tfrac18,\qquad b_{\mathrm{sp}}=\tfrac1{16}.
\]
For a label \(i=(j,\sigma,n)\), \(n=(\bar n,n_3)\in\mathbb Z^3\), set
\begin{equation}
 \mathcal O_i(t)=
 (\widehat E_j^\sigma+\bar n)
 \times[n_3-a-2t-b_{\mathrm{sp}},
        n_3+a+2t+b_{\mathrm{sp}}].
 \label{jet:enlargements}
\end{equation}
These enlargements will only localize upper bounds; they will not
replace the stripe masks when their average is computed.

\begin{lemma}[Disjoint enlargements]
\label{jet:geometry}
For \(0\le t\le\overline T\), the sets \(\mathcal O_i(t)\) are
pairwise disjoint.  Moreover,
\[
 \widehat E_j^\sigma\subset
 \{s_j<|y|<e^\eta s_j\}\subset B(0,s_0).
\]
At each point of the periodic lift there is therefore at most one
label whose enlargement contains that point.
\end{lemma}

\begin{proof}
Opposite stripe patterns at the same scale have distance at least
\(p_j/4\), whereas \(2\delta_j=c_\eta p_j/2<p_j/4\).
The radial gap between the cores for the adjacent scales \(s_j\)
and \(s_{j+1}\) is
\[
 c_\eta(p_j+p_{j+1})
   >\delta_j+\delta_{j+1}.
\]
The gaps between nonadjacent cores are larger.  Enlarging by
\(\delta_j<c_\eta p_j\) still leaves the \(j\)-th set inside its
original open annulus.  This proves disjointness of all the
untranslated transverse enlargements.  They lie in \(B(0,s_0)\),
so their distinct integer translates are disjoint as well.
Finally,
\[
 a+2\overline T+b_{\mathrm{sp}}=\tfrac14.
\]
The axial intervals belonging to distinct integers \(n_3\) are
disjoint.  These two observations prove the assertion for the
products in \eqref{jet:enlargements}.
\end{proof}

Fix a nonnegative smooth probability density \(J\), supported in the
unit ball of \(\R^3\).  Write
\[
 S(\ell)=\operatorname{sech}^2\ell,\qquad
 a_s^\sigma(z)=\sigma+\alpha\tanh(\gamma z/s),
 \qquad d_s=d_0s^{10}.
\]
The positive parameter \(d_0\le1\) will be decreased below finitely
many times.  In particular, we require that
\begin{equation}
 d_0\overline s^{10}\le\tfrac1{64},\qquad
 \overline T d_0\overline s^8\le c_\eta/8 .
 \label{jet:thermal-preliminary}
\end{equation}
The width \(1/64\), including the fixed Gaussian widenings used
below, lies within the range of Lemma~\ref{lem:sheet-bound}.
For \(i=(j,\sigma,n)\), with \(s=s_j\), define on the lift
\[
 c_i^0(x,v)=
 \ind_{E_j^\sigma}(y-\bar n)\ind_{(-a,a)}(z-n_3)
 \frac K s S\!\left(\frac{\gamma(z-n_3)}s\right)d_s^{-3}
 J\!\left(\frac{v-a_s^\sigma(z-n_3)e_3}{d_s}\right).
\]
The periodic datum is
\begin{equation}
 F_0(x,v)=\sum_{j\ge1}\sum_{\sigma=\pm1}
                 \sum_{n\in\mathbb Z^3}c_{(j,\sigma,n)}^0(x,v).
 \label{eq:datum}
\end{equation}
At a given spatial point at most one summand is nonzero.  Integer
translations permute the labels, so the sum is periodic.

\begin{lemma}[Integrability of the datum]
\label{lem:datum}
The datum \eqref{eq:datum} satisfies
\[
 \int_\Omega F_0(1+|v|^2+|\log F_0|)\,\dd x\,\dd v<\infty .
\]
Its truncations \(F_0^{(N)}\), retaining \(j\le N\), converge to
\(F_0\) in mass, energy, and relative entropy.
In addition, if \(\chi\ge0\) is a fixed bounded smooth compactly
supported function of \((x,v)\), periodized after replacing \(x\)
by \(x/s_N\), and \(0\le\lambda_N\le L_0<\infty\), then
\[
 F_0^{(N)}+\lambda_Ns_N\chi(x/s_N,v)
\]
has the same limiting mass, energy, and relative entropy as \(F_0\).
For the last assertion \(s_0\) may be decreased so that distinct
periodic copies of the spatial support of the seed are disjoint.
\end{lemma}

\begin{proof}
Integration over one period and summation over integer copies equal
integration of each unshifted prototype over \(\R^3_x\).
The transverse area of \(E_j^\sigma\) is at most \(C_\eta s_j^2\),
and
\[
 \int_{-a}^a S(\gamma z/s)\,\dd z
 \le \frac s\gamma\int_\R S(\ell)\,\dd\ell
 =\frac{2s}{\gamma}.
\]
The velocity factor integrates to one.  Thus the mass of each
annular sign prototype is at most \(Cs^2\).
Its velocities satisfy
\[
 |v|\le1+\alpha+d_s<2,
\]
which gives the same bound for its energy.

On the set where the density is positive, decompose its logarithm
into the three terms
\[
 \log((K/s)d_s^{-3}),\qquad
 \log S(\gamma z/s),\qquad \log J(w).
\]
The first has absolute value at most \(C(1+|\log s|)\).
The second is integrable against its own weight since
\[
 \int_\R S(\ell)|\log S(\ell)|\,\dd\ell<\infty:
\]
on \(|\ell|\ge1\), its integrand is bounded by
\(C(1+|\ell|)e^{-2|\ell|}\), and on a bounded interval it is bounded.
Also \(\int J|\log J|<\infty\), because \(J\) is bounded and
compactly supported and \(r|\log r|\) is bounded for \(0\le r\le1\).
Consequently the absolute entropy of an annular sign prototype is
at most
\[
 Cs^2(1+|\log s|).
\]
The sums of these bounds over the geometric scales converge.
Since the spatial supports of the prototypes are disjoint modulo
periodicity, the same estimates prove convergence of the truncated
mass, energy, and absolute entropy.  As \(-\log M\) is a constant
plus \(|v|^2/2\), they prove relative entropy convergence as well.

For the seed \(h_N^0=\lambda_Ns_N\chi(x/s_N,v)\), a change of spatial
variables gives
\[
 \int h_N^0(1+|v|^2)\le Cs_N^4,\qquad
 \int h_N^0|\log h_N^0|
       \le Cs_N^4(1+|\log s_N|).
\]
The constants are uniform in \(\lambda_N\in[0,L_0]\), including
\(\lambda_N=0\), since \(\lambda|\log\lambda|\) is bounded on
bounded intervals.
It remains to check overlap with the cold datum.  At every point
\[
 0\le F_0^{(N)}\le C s_N^{-31}.
\]
For \(f,b\ge0\), integration of \(1+\log(f+r)\) over \(0<r<b\)
gives, with the convention at \(b=0\),
\[
 |(f+b)\log(f+b)-f\log f|
 \le b\{2+\log^+(f+b)\}+b|\log b|.
\]
Indeed the positive part of the logarithm is bounded by
\(\log^+(f+b)\), and
\(\int_0^b\log^-(f+r)\,\dd r
 \le\int_0^b\log^-r\,\dd r\le b(1+|\log b|)\).
On the seed support,
\(\log^+(F_0^{(N)}+h_N^0)\le C(1+|\log s_N|)\).
The entropy perturbation is therefore
\(O(s_N^4(1+|\log s_N|))\).  The mass and Maxwellian-reference terms
are controlled by the seed's mass and energy estimates.  This
proves the last assertion.
\end{proof}

\subsection{A transported Maxwellian barrier}

For a fixed scale and sign, the equation
\begin{equation}
 u_s^\sigma(t,z)=a_s^\sigma(z-tu_s^\sigma(t,z))
 \label{jet:scalar-flow}
\end{equation}
defines \(u_s^\sigma\) for all \(t\ge0,z\in\R\).  Indeed
\(\tau\mapsto\tau+t a_s^\sigma(\tau)\) is strictly increasing,
has derivative at least one, and maps \(\R\) onto \(\R\).
In particular, \(|u_s^\sigma|\le1+\alpha\).
For label \(i=(j,\sigma,n)\), abbreviate
\[
 u_i(t,z)=u_{s_j}^\sigma(t,z-n_3),\quad
 \tau_i=z-n_3-tu_i(t,z),\quad w_i=v-u_i(t,z)e_3,
\]
and set
\[
 D_{y,i}=\dist(y-\bar n,E_j^\sigma),\qquad
 D_{z,i}=\dist(\tau_i,[-a,a]).
\]
Choose \(c_{\mathrm{sp}}=1/8\) and a fixed sufficiently large
\(C_m\), depending only on \(J\), and define, for \(t>0\),
\begin{equation}
 m_i(t,x,v)=
 C_m\frac K{s_j}d_{s_j}^{-3} e^t
 \exp\left[-\frac{(2-t)|w_i|^2}{d_{s_j}^2}
 -\frac{c_{\mathrm{sp}}(D_{y,i}^2+D_{z,i}^2)}
              {d_{s_j}^2t}\right].
 \label{jet:majorant}
\end{equation}
Its trace at \(t=0\) is defined by the limit along free
characteristics.  Boundary values of the spatial masks are
irrelevant.  Denote the free transport of the \(i\)-th initial
color by
\[
 f_i(t,x,v)=c_i^0(x-tv,v).
\]

\begin{lemma}[Barrier identities]
\label{jet:barrier}
The constant \(C_m\) can be chosen so that \(f_i\le m_i\) on
\([0,\overline T]\), uniformly in the label.  Moreover, almost
everywhere along every free characteristic,
\begin{equation}
 \D m_i\ge
 \left(1+\tfrac78 |w_i|^2/d_{s_j}^2\right)m_i .
 \label{jet:barrier-derivative}
\end{equation}
There is a constant \(C_\nu\), independent of the label, such that
\begin{equation}
 P(m_i,m_i)=m_i\nu_{m_i},\qquad
 \nu_{m_i}(v)\le C_\nu\frac K{s_j}(d_{s_j}+|w_i|).
 \label{jet:barrier-collision}
\end{equation}
Every free core remains in its own spatial enlargement:
\[
 f_i(t,x,v)>0\quad\Longrightarrow\quad x\in\mathcal O_i(t).
\]
\end{lemma}

\begin{proof}
Suppress the indices, write \(d=d_s\), and put
\(A=(a_s^\sigma)'(\tau)\ge0\).  Implicit differentiation gives
\begin{equation}
 u_z=\frac A{1+tA}\ge0,\qquad
 u_t=-u u_z,\qquad
 \D u=(v_3-u)u_z,\qquad
 \D\tau=\frac{v_3-u}{1+tA}.
 \label{jet:flow-identities}
\end{equation}
Consequently
\[
 \D|w|^2=-2u_z w_3^2.
\]
The distance functions are Lipschitz, hence absolutely continuous
along flights, and at their almost-everywhere differentiability
points
\[
 |\D D_y|\le|\bar w|,\qquad
 |\D D_z|\le |w_3|/(1+tA)\le|w_3|.
\]
Differentiating the logarithm of \eqref{jet:majorant} now gives
\begin{align*}
 \D\log m
 ={}&1+\frac{|w|^2}{d^2}
      +\frac{2(2-t)u_z w_3^2}{d^2}\\
    &+\frac{c_{\mathrm{sp}}(D_y^2+D_z^2)}{d^2t^2}
      -\frac{2c_{\mathrm{sp}}(D_y\D D_y+D_z\D D_z)}{d^2t}.
\end{align*}
For each of the two spatial terms use
\(r^2-2r q\ge-q^2\).  This yields
\[
 \D\log m\ge
 1+(1-c_{\mathrm{sp}})|w|^2/d^2
   +2(2-t)u_z w_3^2/d^2.
\]
The last term is nonnegative on the preliminary time interval,
proving \eqref{jet:barrier-derivative}.  In particular the large
derivative of the initial axial profile causes no adverse
\(s^{-1}\) term.

To compare with free transport, start a core flight at
\(x^0=(y^0,z^0)\) in its initial support and write
\[
 v=a_s^\sigma(z^0-n_3)e_3+d w^0,\qquad |w^0|\le1.
\]
At \(x=x^0+tv\), the function
\(p\mapsto p-a_s^\sigma(z-n_3-tp)\) has derivative at least one,
its value at \(v_3\) is \(dw_3^0\), and its value at \(u\) is zero.
Thus
\[
 |v_3-u|\le d|w_3^0|,\quad |\bar v|=d|\bar w^0|,\quad
 D_y\le td|\bar w^0|,\quad D_z\le td|w_3^0|.
\]
In particular \(|w|\le d\) and
\((D_y^2+D_z^2)/(d^2t)\le t\).
The exponential in \eqref{jet:majorant}, including \(e^t\), is
bounded below on this flight by \(e^{-3}\).
Taking \(C_m\ge e^3\|J\|_\infty\) gives \(f_i\le m_i\), since
\(S\le1\).
The same calculation supplies the required initial domination.
Away from the closed initial spatial support the trace of \(m_i\)
is zero; on it the distance exponent tends to zero along flights.
The inequalities therefore integrate from time zero.  No
differentiation of a spatial indicator is used.

At a fixed point \((t,x)\), \(m_i\) is a scalar multiple of the
Maxwellian with center \(u_i e_3\) and exponent
\((2-t)/d_s^2\).  Conservation of momentum and energy in an elastic
collision gives \(m_i'm_{i,*}'=m_i m_{i,*}\), which proves the
first identity in \eqref{jet:barrier-collision}.  Its velocity mass
is at most \(CK/s\), and its first moment about \(u_i e_3\) is at
most \(CKd_s/s\), by the substitution
\(p=u_i e_3+d_s\xi\).  The inequality
\(|v-p|\le|w_i|+|p-u_i e_3|\) proves the frequency bound.

Finally, the free-flight estimates and
\eqref{jet:thermal-preliminary} give
\[
 D_y\le\overline T d_s\le\delta_j/2.
\]
Its axial velocity has modulus less than \(2\), so its position
lies between \(n_3-a-2t\) and \(n_3+a+2t\).
It belongs to \(\mathcal O_i(t)\), as asserted.
\end{proof}

\subsection{Cold control and the density Jacobian}

The labels for the finite approximation are
\[
 \mathcal I_N=\{(j,\sigma,n):1\le j\le N,\ \sigma=\pm1,\
                                  n\in\mathbb Z^3\}.
\]
For \(i=(j,\sigma,n)\) we also write \(s_i=s_j\).
The infinitely many integer-copy labels merely implement
periodicity: translating \(x\) permutes them.  We consider
nonnegative colored mild solutions with these initial colors and
any nonnegative hot component, with a common finite collision
cutoff \(b>0\).  The statement below applies throughout every
interval on which the finite-approximation solution exists with
bounded weighted sup norm.  Its constants do not use that norm or
any bound on the hot component.

\begin{proposition}[Uniform cold control]
\label{prop:cold-control}
There is a constant \(L\), depending on the fixed parameters, such
that after imposing the additional condition
\begin{equation}
 Ld_0\overline s^9\le1,
 \label{jet:thermal-comparison}
\end{equation}
the following statements hold at every time of existence of the
finite approximation in \(0\le t\le\overline T\).
\begin{enumerate}[label=\textup{(\roman*)}]
 \item For every retained label \(i=(j,\sigma,n)\),
 \begin{equation}
  0\le c_i\le f_i+\varepsilon_j m_i,\qquad
  \varepsilon_j=Ld_{s_j}/s_j .
  \label{jet:cold-remainder}
 \end{equation}
 \item At each spatial point call the unique label with
 \(x\in\mathcal O_i(t)\), if one exists and is retained, the main
 label.  Let \(c_{\mathrm{main}}\) be that color, or zero, and set
 \(e=c-c_{\mathrm{main}}\).  For every \(A\ge0\),
 \begin{equation}
  0\le e(t,x,v)\le C_A t^A e^{-4|v|^2}.
  \label{jet:escaped}
 \end{equation}
 The upper bound remains valid after restricting the escaped sum
 to any subset of its labels.
 \item If the main label has scale \(s\), its velocity measure is
 bounded by a fixed multiple of \((t+s)^{-1}\) times a sum of
 admissible probability sheets from Lemma~\ref{lem:sheet-bound}.
 Its free-core contribution has the exact parameterization
 \eqref{eq:core-jacobian} below.  The remainder has the smaller
 velocity-mass coefficient \(Cd_s/s^2\).
 In particular, for every fixed \(k\ge0\),
 \begin{equation}
  \int_{\R^3}\langle v\rangle^k c_{\mathrm{main}}\,\dd v
       \le\frac{C_k}{t+s},\qquad
  \int_{|v|>R}\langle v\rangle^k c_{\mathrm{main}}\,\dd v
       \le\frac{C_k e^{-c_kR^2}}{t+s}\quad(R\ge3).
  \label{jet:moment-envelope}
 \end{equation}
 \item For the mixed-cold source in the hot equation,
 \begin{equation}
  \sum_{\substack{i,j\in\mathcal I_N\\i\ne j}}
       P^b(c_i,c_j)(t,x,v)
       \le C_A t^{A-1}H(v)\qquad (t>0)
  \label{jet:mixed-source}
 \end{equation}
 for every \(A\ge0\), with constants uniform in \(N\) and \(b\).
\end{enumerate}
\end{proposition}

\begin{proof}
\emph{Nonlinear comparison.}
For one color put \(U_i=f_i+\varepsilon_jm_i\).
Lemma~\ref{jet:barrier} gives \(U_i\le(1+\varepsilon_j)m_i\).
By \eqref{jet:thermal-comparison}, \(\varepsilon_j\le1\); therefore
\[
 P^b(U_i,U_i)\le P(U_i,U_i)
 \le4 C_\nu\frac K{s_j}(d_{s_j}+|w_i|)m_i.
\]
On the other hand,
\[
 \D U_i\ge
 L\frac{d_{s_j}}{s_j}
       (1+\tfrac78|w_i|^2/d_{s_j}^2)m_i .
\]
The ratio
\((1+r)/(1+\tfrac78r^2)\) is bounded on \(r\ge0\).
Choose \(L\ge1\) at least \(4C_\nu K\) times that bound.
Then \(\D U_i\ge P^b(U_i,U_i)\).
This choice is made before decreasing \(d_0\), so
\eqref{jet:thermal-comparison} is compatible with all earlier
constraints.

The cold equation implies \(\D c_i\le P^b(c_i,c_i)\), because its
loss is nonnegative.  To justify comparison explicitly, integrate
both inequalities along a free path, using
\(c_i(0)\le U_i(0)\), and put \(z_i=(c_i-U_i)_+\).
Positivity and bilinearity give
\[
 z_i(t)\le\int_0^t
  \bigl[2P^b(U_i,z_i)+P^b(z_i,z_i)\bigr]
       (\tau,x-(t-\tau)v,v)\,\dd\tau .
\]
On any interval of bounded weighted norm, the gain bound
Lemma~\ref{lem:hot-bound} implies
\[
 \|z_i(t)/H\|_\infty
 \le C_I\int_0^t\|z_i(\tau)/H\|_\infty\,\dd\tau .
\]
Indeed \(m_i/H\) is bounded for a fixed scale, uniformly on the
preliminary interval, and
\(z_i\le c_i+U_i\) supplies the other bounded factor.
Iteration of this integral inequality, or its scalar integrating
factor, gives \(z_i=0\).  This proves
\eqref{jet:cold-remainder} without a hot bound.
The argument is unchanged for each periodic-copy label.

\emph{Summation of escaped tails.}
For a nonmain label, \eqref{jet:enlargements} and
\(|u_i|\le1+\alpha<2\) imply either \(D_{y,i}\ge\delta_j\) or
\(D_{z,i}\ge b_{\mathrm{sp}}\).
Since \(\delta_j\le b_{\mathrm{sp}}\), in either case
\[
 D_{y,i}^2+D_{z,i}^2\ge
       (c_\eta/4)^2s_j^4.
\]
If \(a_v=(2-t)/d_s^2\), then \(a_v\ge8\), and completing the
velocity square gives
\[
 \sup_v \exp\{4|v|^2-a_v|v-ue_3|^2\}
   =\exp\!\left(\frac{4a_v|u|^2}{a_v-4}\right)
   \le e^{8(1+\alpha)^2}.
\]
Thus \(c_i\le2m_i\) bounds its escaped contribution by
\begin{equation}
 C s_j^{-31}e^{-4|v|^2}
 \exp\!\left[-\frac{c_{\mathrm{sp}}
                     (D_{y,i}^2+D_{z,i}^2)}{d_{s_j}^2t}\right].
 \label{jet:tail-prefactor}
\end{equation}
The factor \(d_0^{-3}\) has been included in the fixed constant.

For completeness the integer-copy sum is uniform in the scale.
Since \(E_j^\sigma\subset[-s_0,s_0]^2\),
\[
 D_{y,i}^2\ge
   \sum_{k=1}^2(|y_k-n_k|-s_0)_+^2,\qquad
 D_{z,i}\ge
   (|z-n_3|-a-2\overline T)_+.
\]
For \(\beta>0,R\ge0,\epsilon>0\), comparison of the two monotone
Gaussian tails with their integrals gives
\[
 \sup_{\xi\in\R}\sum_{n\in\mathbb Z}
 e^{-\beta(|\xi-n|-R)_+^2/\epsilon^2}
 \le 2R+5+\epsilon\sqrt{\pi/\beta}.
\]
There are at most \(2R+2\) integers in the central interval; for
each of the remaining tails, keep its first term and bound the
rest by the corresponding Gaussian integral, which proves the
displayed bound.  Applying this estimate to the three coordinates, with
\(\epsilon=d_s\sqrt t\), bounds the lattice sum of one half of the
spatial exponent uniformly for \(t\le\overline T\).
Use the other half to retain the escaped distance.  For some fixed
\(\beta_0>0\) this gives
\[
 \sum_{n:\,i\ {\rm nonmain}}
 e^{-c_{\mathrm{sp}}(D_{y,i}^2+D_{z,i}^2)/(d_s^2t)}
 \le C\exp[-\beta_0/(s^{16}t)] .
\]
Here \(\beta_0\) includes \(d_0^{-2}\), and the two signs only
contribute a factor two.

For arbitrary \(A\ge0\) and \(B>0\), choose \(p\ge A\) with
\(16p-31\ge B\).  From \(e^{-x}\le(p/e)^p x^{-p}\), and \(t\le1\),
\[
 s^{-31}e^{-\beta_0/(s^{16}t)}
       \le C_{A,B}t^As^B .
\]
The geometric sum of \(s_j^B\) is finite.  More precisely, restricting
to \(s_j\le\rho\) gives a bound \(C_{A,B}t^A\rho^B\).
This proves \eqref{jet:escaped}; it also shows that any fixed
extra power of \(s_j^{-1}\) can be absorbed by increasing \(p\).
All these estimates were obtained by summing upper bounds over
every scale and copy, so they are uniform in \(N\).

\emph{The exact free-core velocity measure.}
Fix \(t,x\) and one label, abbreviating
\(z_c=z-n_3\), \(y_c=y-\bar n\), \(s=s_j\), \(d=d_s\).
For \(w\in\R^3\) define \(\Phi(w_3)\) by
\[
 \Phi(w_3)=a_s^\sigma(z_c-t\Phi(w_3))+dw_3,\qquad
 \bar p=d\bar w,\qquad p_3=\Phi(w_3).
\]
The same strictly increasing map used in
\eqref{jet:scalar-flow} gives a unique \(\Phi\), and
\[
 \Phi'(w_3)=
 \frac d{1+t\alpha\gamma
           S(\gamma(z_c-t\Phi(w_3))/s)/s}.
\]
The full velocity change of variables therefore gives exactly
\begin{equation}
 \begin{aligned}
 f_i(t,x,p)\,\dd p
 ={}&
 \frac{K S(\gamma(z_c-tp_3)/s)}
      {s+t\alpha\gamma S(\gamma(z_c-tp_3)/s)}
 J(w)\\[-2pt]
 &{}\times
 \ind_{E_j^\sigma}(y_c-td\bar w)
 \ind_{(-a,a)}(z_c-tp_3)\,\dd w .
 \end{aligned}
 \label{eq:core-jacobian}
\end{equation}
For \(k_0=\alpha\gamma>0\) and \(0<S\le1\),
\begin{equation}
 \frac{KS}{s+t k_0S}
 \le\frac K{s+t k_0}
 \le\frac{K\max(1,k_0^{-1})}{s+t}.
 \label{jet:density-coefficient}
\end{equation}
Thus this estimate is uniform even when \(t/s\) is large or
the value of \(S\) is small.

We spell out the product-sheet upper bound, since the axial law
can be very compressed.  Let
\[
 g_2(\bar w)=\pi^{-1}e^{-|\bar w|^2},\qquad
 g_1(w_3)=\pi^{-1/2}e^{-w_3^2}.
\]
There is a fixed \(C_J\) such that \(J(w)\le C_Jg_2(\bar w)g_1(w_3)\).
The axial map \(\Phi\) depends only on \(w_3\), and
\[
 |\Phi(w_3)-u_s^\sigma(t,z_c)|\le d|w_3|
\]
by the derivative-at-least-one argument.
The pushforward of \(g_2(\bar w)\dd\bar w\) by
\(\bar w\mapsto d\bar w\) is a centered transverse Gaussian.
Independently, let \(\rho\) be the pushforward of
\(g_1(w_3)\dd w_3\) by \(\Phi\).
It is a probability law and, uniformly in \(s,t,x\),
\[
 \int_\R e^{8p_3^2}\,\rho(\dd p_3)
 \le e^{16(1+\alpha)^2}
       \int_\R e^{16d^2w_3^2}g_1(w_3)\,\dd w_3
 \le C.
\]
Equations \eqref{eq:core-jacobian} and
\eqref{jet:density-coefficient}, after discarding the masks,
bound the core measure by \(C/(t+s)\) times this product law.
No density bound on \(\rho\) is asserted or needed.

For the remainder, use \(p=u_i e_3+d_s w\) in
\eqref{jet:majorant}.  Its measure is bounded by
\[
 \varepsilon_jm_i\,\dd p
 \le C\frac{d_s}{s^2}e^{-(2-t)|w|^2}\,\dd w,
\]
which is again a product-sheet majorant, with the stated smaller
coefficient.  Since
\[
 (d_s/s^2)(t+s)=d_0s^8(t+s)
       \le d_0\overline s^8(\overline T+\overline s),
\]
this coefficient is also at most \(C/(t+s)\).
The width restrictions in \eqref{jet:thermal-preliminary} allow
these Gaussian laws and their fixed, slightly wider majorants
in Lemma~\ref{lem:sheet-bound}. Lemma~\ref{ker:parameter-sheets}
applies with the common axial bound \(A=1+\alpha\), the fixed
Gaussian majorant constant, and the displayed physical widths.
The factors \(C/(t+s)\) and \(Cd_s/s^2\) remain outside the
probability-sheet normalization. A parameter restriction \(|w|>R\)
therefore contributes at most \(Ce^{-cR^2}\) times its respective
amplitude to the normalized gain. For the remainder,
\(2-t\ge2-\overline T>1\), so its parameter Gaussian is bounded
by \(e^{-|w|^2}\). The normalized kernel-tail
conclusions of Lemma~\ref{lem:kernel-tightness} also apply, with
these amplitudes.

The core velocities are supported in \(|p|<2\).
The remainder has uniformly small Gaussian width and a center in
\(|p_3|\le1+\alpha\).  Substitution in its Gaussian integral gives
all its polynomial moments and the uniform Gaussian moment
tails.  Combined with \eqref{jet:density-coefficient}, these give
\eqref{jet:moment-envelope}.

\emph{The mixed source.}
At a fixed point, separate the main color from \(e\).  Positivity
and the ordered nature of the cross-color sum give
\[
 \sum_{i\ne j}P^b(c_i,c_j)
       \le 2P(c_{\mathrm{main}},e)+P(e,e).
\]
The function \(e^{-4|v|^2}/H(v)\) is essentially bounded, so
\eqref{jet:escaped} bounds \(e\) by \(C_A t^AH\).
Lemma~\ref{lem:sheet-bound} and the product measures just proved
bound \(P(c_{\mathrm{main}},H)\) by \(C(t+s)^{-1}H\);
when there is no main color this term is zero.
Lemma~\ref{lem:hot-bound} bounds \(P(H,H)\) by \(CH\).
Since \(t+s\ge t\) and \(t\le1\),
\[
 2P(c_{\mathrm{main}},e)+P(e,e)
 \le C_A\{t^A/(t+s)+t^{2A}\}H
 \le C_A' t^{A-1}H .
\]
This proves \eqref{jet:mixed-source}.
\end{proof}

It is useful to record the spatial form of the moment bound.
For a fixed \(k\ge0\), Proposition~\ref{prop:cold-control} implies
\begin{equation}
 \int_{\R^3}\langle v\rangle^k c(t,x,v)\,\dd v
 \le C_k\sum_{i\in\mathcal I_N}
       \frac{\ind_{\mathcal O_i(t)}(x)}{t+s_i}
       +C_{A,k}t^A.
 \label{jet:spatial-moment-envelope}
\end{equation}
At most one term of the displayed sum is present at any point.
In one periodic cell, the spatial volume of the enlargement
belonging to a fixed annular sign, with its periodic copies
included, is at most \(Cs_i^2\).
The same bound with a factor tending to zero as the velocity
cutoff increases holds for the integrated velocity tails.

\subsection{Quantitative survival on a free core}

The upper estimates did not use a hot bound.  The following
consequence, used only for the dilation, adds the cap assumption.

\begin{corollary}[Survival of a cold core under the cap]
\label{jet:survival}
Suppose a finite approximation satisfies
\[
 0\le h(t,x,v)\le(t+s_N)H(v)
\]
on an initial interval \([0,T]\subset[0,\overline T]\).
All conclusions below concern times \(t\in[0,T]\).
For every \(M\ge0\) there are constants \(C,C_M\), independent of
\(T,N\), the cutoff, and the seed, such that
\begin{equation}
\begin{gathered}
 e^{-E_{s,N}(t)}f_i(t,x,v)
       \le c_i(t,x,v)
       \le f_i(t,x,v)+L(d_s/s)m_i(t,x,v),\\
 E_{s,N}(t)=C(d_st/s+t^2+s_Nt)+C_Mt^{M+1}
\end{gathered}
 \label{jet:survival-sandwich}
\end{equation}
for every label of scale \(s\).
Consequently, if \(r\to0\), \(s_N\le r\), \(s/r\in[\epsilon,R]\)
for fixed \(0<\epsilon<R<\infty\), and \(0\le t\le\min(T,Ar)\), then
\[
 E_{s,N}(t)\le
 C_{\epsilon,R,A}(d_s+r^2)+C_{M,A}r^{M+1}\longrightarrow0
\]
uniformly in the label and in points of its free core.
The transverse free displacement divided by the stripe period
satisfies
\begin{equation}
 \frac{td_s}{s^2}\le
       C_{\epsilon,R,A}r^9\longrightarrow0.
 \label{jet:stripe-displacement}
\end{equation}
On these scales the difference between cold and free-core
velocity measures also satisfies, uniformly in position and time,
\begin{equation}
 r\int_{\R^3}|c_i-f_i|(t,x,v)\,\dd v\longrightarrow0.
 \label{jet:survival-measure}
\end{equation}
\end{corollary}

\begin{proof}
The loss-integrating-factor formula for the \(i\)-th mild equation
and nonnegativity of its gain give
\[
 c_i(t,x,v)\ge f_i(t,x,v)
 \exp\left[-\int_0^t
   \nu_F^b(\tau,x-(t-\tau)v,v)\,\dd\tau\right].
\]
Only points with \(f_i(t,x,v)>0\) need be considered.
Their entire backward flight, equivalently the corresponding
forward flight from the initial support, is in
\(\mathcal O_i(\tau)\) by Lemma~\ref{jet:barrier}.
Also \(|v-u_i(\tau,z(\tau))e_3|\le d_s\) throughout it.

At such a position the free part of the same color has velocities
within \(d_s\) of this center, by
\eqref{eq:core-jacobian}, and mass at most \(C/(\tau+s)\).
Its contribution to the frequency at \(v\) is at most
\(Cd_s/(\tau+s)\le Cd_s/s\).
The remainder has mass \(Cd_s/s^2\) and first moment about the
center at most \(Cd_s^2/s^2\).  Since \(d_s/s\le1/L\), its
contribution is at most \(Cd_s/s\) as well.
The same-color frequency is therefore bounded by \(Cd_s/s\).

Every other color is escaped along this flight.  Its total
frequency is at most \(C_M\tau^M\), by
\eqref{jet:escaped} and \(|v|<2\).
The cap applies at every preceding time \(\tau\in[0,t]\), so
\[
 \nu_h(v)\le 2\pi(\tau+s_N)
       \int_{\R^3}(|v|+|p|)H(p)\,\dd p
       \le C(\tau+s_N)
\]
at these bounded velocities.  Since \(\nu_F^b\le\nu_F\),
integration proves the optical-depth bound in
\eqref{jet:survival-sandwich}.  Its upper bound is
\eqref{jet:cold-remainder}.

If \(s/r\in[\epsilon,R]\), \(t\le Ar\), and \(s_N\le r\), the
first term in the optical depth is bounded by
\(C_{A,\epsilon}d_s\), and the two hot terms are \(O_A(r^2)\).
All stated convergence is uniform.  Moreover
\[
 td_s/s^2=d_0t s^8\le d_0 A R^8r^9,
\]
which proves \eqref{jet:stripe-displacement}.

For the last claim the sandwich gives
\[
 |c_i-f_i|\le
       (1-e^{-E_{s,N}(t)})f_i+L(d_s/s)m_i.
\]
Integrating in velocity and using the core and remainder density
bounds yields
\[
 r\int|c_i-f_i|\,\dd v
 \le C E_{s,N}(t)\frac r{t+s}
       +C\frac{r d_s}{s^2}.
\]
Here \(r/(t+s)\le1/\epsilon\), and
\(r d_s/s^2=d_0r s^8\le C_Rr^9\).
Both terms tend to zero uniformly, proving
\eqref{jet:survival-measure}.
\end{proof}
